\section{A certain liability need not identify its creditor}
\label{sec:query-certainty}

Suppose the admitted records establish exactly one payable of 100 currency
units, but leave its creditor unresolved between A and B. For this constructed
example the declared family has two histories: in $h_A$ the company owes A
100 and owes B nothing; in $h_B$ it owes B 100 and owes A nothing. These are
complete alternatives within the example, not two sampled points from a wider
family. The question \emph{does a payable exist?} has the answer yes in both
histories. The question \emph{is A owed 100?} does not.

Now ask for the rows of the payable schedule. If $Q(h)$ returns those rows,
then $Q(h_A)=\{(A,100)\}$ and $Q(h_B)=\{(B,100)\}$. Their intersection is
empty: no creditor--amount row occurs in both. Reading that empty table as
the company's complete schedule would wrongly assert that it owes nothing.
Reading the union as its complete schedule would instead assert that it owes
200. Evaluating the total within each history gives 100 in both. The
uncertainty concerns the creditor, not the existence or total of the payable.

This distinction is central to Libkin's critique of intersection-based
answers. His example retains a tuple with an unknown second component;
discarding it loses knowledge that a tuple with the stated first component
exists. An intersection read under a closed-world interpretation can also
add an unsupported absence claim \citep[Section~3 and the opening of
Section~6]{libkin2014}. A table of individually certain rows is therefore a
restricted view of knowledge, not generally a replacement for the possible
answer schedules.

For a nonempty feasible history family $H=\Hist_M(D_k)$ and a total,
deterministic query $Q$ on its histories, retain the family of answers
\[
 \mathcal{A}_Q(H)=\{Q(h):h\in H\}.
\]
Each answer remains linked to the histories that produce it. Equality of
payable schedules alone does not license merging histories for another query.
For a row-valued query, a separately requested certain-tuples view is
\[
 T_Q(H)=\bigcap_{h\in H}Q(h).
\]
Membership in $T_Q(H)$ asserts that a row appears in every answer. Omission
asserts only that the row is not established in every answer; it does not
assert absence from every answer. A total computed after taking this
intersection answers a different question from totals computed within
histories. In the payable example, $T_Q(H)$ is empty while
$\{\operatorname{total}(Q(h)):h\in H\}=\{100\}$.

For a proposition $P$ about a concrete answer, Boolean certainty has the
separate meaning
\[
 \operatorname{Certain}_H(P\circ Q)
 \quad\Longleftrightarrow\quad
 \forall h\in H,\ P(Q(h)).
\]
Taking $P$ to mean that the schedule contains at least one payable proves
existence without identifying A or B as the creditor. Allowing a third history
with no payable would change that answer. The two-history result consequently
depends on the admitted evidence and scope that require the liability to
exist, not merely on the two rows chosen for display.

Libkin distinguishes a logical description of knowledge shared by all
answers from an object representing their greatest lower bound in an
information ordering. These representations need not exist in the chosen
language; ordinary tuple intersection is not their general definition
\citep[Section~5.3, equations~(6)--(8)]{libkin2014}. AccDSL's answer family and
Boolean propositions above are local definitions. They do not assert that
every family has a finite relational representation. A symbolic answer must
carry its conditions and preserve the stated family, or declare the
restriction or approximation it makes.

Input completeness and output meaning must also be explicit. A model that
fixes the available event templates and contract slots imposes a scope
restriction; it does not establish that a public filing lists every economic
event. Nor does that restriction entitle a display to interpret every omitted
row as false. The naive-evaluation results reported by Libkin require suitable input and
output semantics, monotonicity in their information order, genericity, and an
answer representation system. His relational examples cover positive queries
and a specified guarded extension under closed-world semantics
\citep[Sections~6.1--6.2, equations~(9)--(10)]{libkin2014}. Those results do not
authorize evaluating arbitrary financial negation, aggregates or nonlinear
constraints by treating unknown amounts as ordinary constants.

The query contract records whether it returns a proposition, individually
certain rows, or a family of answers, together with the family scope and
completeness and the interpretation of omitted rows. Unsupported evaluation
and unfinished search keep their existing failure or unresolved status. In
particular, a sampled answer family cannot be relabelled as the full family
merely because its displayed rows agree.
